\section{科研智能体系统设计（证据层为核心）}

\begin{designbox}{产品定位}
科研智能体不是一个普通聊天框, 也不只是一套 RAG 检索器。它的差异化在于\textbf{证据立场判定}:
对一句论断, 系统能说出文献是"支持/反驳/中立", 指出证据落在哪个句子, 低置信时如实说
"证据不足", 并把检测到的矛盾沉淀为可读取的\textbf{科学争议地图}。这个原语同时是核查能力、
争议检测器和可被其他 agent 调用的服务。
\end{designbox}

\subsection{核心能力矩阵}

\begin{center}
\begin{tabularx}{\textwidth}{>{\bfseries}p{30mm}X X}
\toprule
能力 & 用户体验 & 后端能力 \\
\midrule
证据问答 & 自然语言问方向, 系统返回答案与证据卡 & 混合检索、GraphRAG、引用约束生成 \\
\multicolumn{1}{l}{\textbf{论断核查}} & \multicolumn{1}{l}{一句话得到支持/反驳/中立, 证据句可展开, 边界诚实} & \multicolumn{1}{l}{stance L3: 句级证据, 类别概率, 校准拒答}\\
争议地图 & 看到"文献里有人反对"的论断与双方证据 & claim\_evidence\_stance + contradictions 视图 + API/Agent tool/MCP resource \\
趋势预测 & 询问未来方向, 返回解释与不确定性 & 关键词时序、增长率、显著性核实 \\
论文推荐 & 输入兴趣, 得到推荐与解释原因 & 向量/主题/作者/时间加权, 可解释因子 \\
图谱查询 & 查询作者/关键词/主题关系 & 图谱索引与 NetworkX 预计算 \\
\bottomrule
\end{tabularx}
\end{center}

\subsection{Stance L3 管线（技术主叙事）}

\begin{tikzpicture}[
  node distance=6mm,
  every node/.style={font=\small},
  box/.style={draw=black, fill=white, rounded corners=2pt, align=center, minimum width=30mm, minimum height=10mm},
  arrow/.style={-{Latex[length=2mm]}, thick, color=black}
]
  \node[box, very thick] (q) {论断 Claim};
  \node[box, right=6mm of q] (ret) {检索相关文献\\6 篇候选};
  \node[box, right=6mm of ret] (sent) {证据句定位\\句级跨度};
  \node[box, below=7mm of sent] (judge) {单条立场判定\\SUPPORT / CONTRADICT / NEUTRAL};
  \node[box, left=6mm of judge] (agg) {聚合 verdict\\强支持/部分/争议/反驳/不足};
  \node[box, left=6mm of agg] (store) {落库\\claim\_evidence\_stance};
  \node[box, below=7mm of store] (map) {争议地图\\contradictions 视图};
  \node[box, densely dashed, above=7mm of judge] (rule) {回退纪律\\LLM 不可用 $\rightarrow$ similarity 仅判证据不足};
  \draw[arrow] (q) -- (ret);
  \draw[arrow] (ret) -- (sent);
  \draw[arrow] (sent) -- (judge);
  \draw[arrow] (judge) -- (agg);
  \draw[arrow] (agg) -- (store);
  \draw[arrow] (store) -- (map);
  \draw[arrow,densely dashed] (rule) -- (judge);
\end{tikzpicture}

\begin{plainbox}
\begin{itemize}
  \item \textbf{句级证据}: 采用的证据能落到句级（或等价最小跨度）, 不拿整段摘要充当证据。
  \item \textbf{校准与拒答}: 低置信倾向"证据不足"; 给出类别概率或分数, 呈校准/拒答表现。
  \item \textbf{限定条件}: 人群/剂量/时间/物种等明显冲突时降级为 NEUTRAL 或证据不足并写理由。
  \item \textbf{对称性}: 论断 C 与反论断 $\neg$C 不得无理由同判为支持（这是早期"相似即支持"的裂缝）。
  \item \textbf{回退纪律}: LLM 不可用时仅按相似度给"证据不足", \textbf{禁止}回退序列输出"强支持"式措辞。
  \item \textbf{金标准}: 双语主集 + SciFact 协议锚点, 与纯相似度 baseline 同集对照, 结果进项目报告。
\end{itemize}
\end{plainbox}

\subsection{协议与模块边界}

\begin{center}
\begin{tabularx}{\textwidth}{p{34mm}X X}
\toprule
接口或模块 & 输入 & 输出 \\
\midrule
\texttt{/api/agent/stream} & 问题、历史、session\_id、retry & SSE 事件流（plan/tool/reflect/final）+ meta \\
\texttt{/api/chat} & 问题、证据问答 & 答案、证据、引用、置信说明 \\
\texttt{/api/search} & 查询词 & 混合检索结果（FTS+向量+RRF） \\
\texttt{/api/trends} & 关键词、时间范围 & 趋势曲线、显著性、解释 \\
\texttt{/api/recommend} & 论文 id / 主题 & 推荐列表、解释因子 \\
\texttt{/api/graph} & 作者/关键词/主题 & 节点、边、社区、中心性 \\
\texttt{verify\_claim (tool)} & 论断 & verdict、证据句、立场统计 \\
\texttt{disputes (MCP resource)} & — & 当前争议前线论断 \\
\bottomrule
\end{tabularx}
\end{center}

\subsection{可信答案原则}

\begin{itemize}
  \item 答案必须绑定证据卡, 证据卡必须能追溯到论文标题、年份、来源与片段。
  \item 核查类结论区分"相关"与"支持"; 拒不出现"证据不足却强说支持"。
  \item 系统敢说"证据不足 / 文献里有人反对", 也接受人类反驳并认错。
  \item 趋势/推荐给明数据显示与不确定性, Web 不把代理指标当用户价值。
\end{itemize}

\clearpage